\chapter{État de l'art et choix des composants}\label{cap:etat-art}

\section{Définir l'approche agentique}

Dans ce mémoire, le terme « agentique » désigne l'organisation de tâches autour d'un état partagé, d'outils et de règles de passage entre étapes. Le prototype initial en retient une forme limitée : LangGraph coordonne les contrôles et les traitements déterministes, puis peut déclencher une synthèse DeepSeek. Il ne s'agit pas d'une équipe autonome d'agents financiers.

J'ai retenu de ReAct et de Toolformer l'usage d'outils externes pour fournir au modèle les résultats dont il a besoin \citep{yao2023react,schick2023toolformer}. Dans mon cas d'usage, les rendements, les coûts et les autres valeurs financières sont calculés en amont, puis transmis au modèle. Le RAG pourrait compléter cette architecture pour traiter des documents ou des actualités financières \citep{lewis2020rag}, mais il n'est pas utilisé dans l'expérience.

Reflexion et AutoGen apportent des pistes pour l'architecture cible, avec le retour critique ou les échanges entre agents \citep{shinn2023reflexion,wu2023autogen}. Ces mécanismes dépassent le besoin de l'expérience initiale de synthèse : le LLM y restitue un état quantitatif déjà produit. Il ne calcule pas le signal, dont la robustesse doit être évaluée séparément de la qualité de la synthèse.


\section{Du besoin du prototype au choix des outils}

La comparaison de LangGraph et de Lean m'a aidé à séparer deux besoins : organiser les étapes d'un workflow et simuler un portefeuille \citep{githubLangGraph,githubLean}. J'ai conservé un moteur Python pour les calculs de quatre ETF, puis LangGraph pour suivre l'état et bloquer la synthèse si le contrôle qualité échoue. Les autres frameworks ont été comparés sur documents, sans essai d'exécution permettant de les classer en performance.

Le tableau~\ref{tab:frameworks-agents} situe les six frameworks examinés. Le choix de LangGraph repose sur les transitions et le suivi de l'état mobilisés dans mon prototype ; les calculs et les règles métier sont des fonctions séparées.

\begin{table}[H]
\centering
\caption{Frameworks agentiques généralistes disponibles sur GitHub.}
\label{tab:frameworks-agents}
\scriptsize
\begin{adjustbox}{width=\textwidth}
\begin{tabular}{p{0.18\textwidth}p{0.24\textwidth}p{0.30\textwidth}p{0.24\textwidth}}
\toprule
\textbf{Projet GitHub} & \textbf{Positionnement} & \textbf{Intérêt pour un desk multifactoriel} & \textbf{Limites principales} \\
\midrule
LangGraph \citep{githubLangGraph} & Workflows d'agents sous forme de graphes, avec état persistant et contrôle fin des transitions. & Adapté aux processus audités : ingestion, calcul de facteurs, revue de risque, validation, journalisation. & Nécessite de développer les outils quantitatifs et les contrôles métier autour du graphe. \\
AutoGen \citep{githubAutoGen} & Conversations multi-agents et orchestration d'équipes d'agents. & Utile pour comparer des rôles de recherche, analyste, contrôleur de risque et superviseur. & Les conversations libres doivent être fortement contraintes pour un usage régulé. \\
CrewAI \citep{githubCrewAI} & Agents orientés rôles, tâches et processus. & Simple pour construire une équipe d'agents de recherche de marché et de reporting. & Moins adapté seul aux workflows critiques nécessitant des états vérifiables et des garde-fous stricts. \\
Semantic Kernel \citep{githubSemanticKernel} & SDK Microsoft pour connecter modèles, plugins, fonctions et mémoire. & Intéressant dans un environnement d'entreprise déjà Microsoft, avec intégration possible à des services internes. & L'architecture de trading reste à construire ; le framework ne fournit pas de modèle financier. \\
CAMEL \citep{githubCamel} & Environnements et méthodes pour sociétés multi-agents. & Pertinent pour expérimenter des interactions entre analystes spécialisés. & Plus orienté recherche agentique que production financière. \\
OpenAI Swarm \citep{githubSwarm} & Exemple éducatif de routage entre agents. & Montre clairement la logique de transfert entre rôles. & Projet expérimental, insuffisant pour une architecture de production. \\
\bottomrule
\end{tabular}
\end{adjustbox}
\end{table}

\section{Socle quantitatif et accès aux données}

Une recherche sur actions individuelles demanderait de définir l'univers, préparer les données, calculer les signaux et construire le portefeuille sous contraintes. J'ai réduit ce travail à VLUE, MTUM, QUAL et USMV, initialement équipondérés et rééquilibrés chaque mois. Les ETF représentent les expositions ; le prototype ne reconstruit pas les facteurs titre par titre. L'extension d'allocation conserve cet univers et compare les propositions DeepSeek à deux règles simples.

\begin{table}[H]
\centering
\caption{Briques quantitatives GitHub utiles pour une étude multifactorielle.}
\label{tab:github-quant}
\scriptsize
\begin{adjustbox}{width=\textwidth}
\begin{tabular}{p{0.18\textwidth}p{0.25\textwidth}p{0.32\textwidth}p{0.21\textwidth}}
\toprule
\textbf{Projet GitHub} & \textbf{Fonction principale} & \textbf{Apport dans l'architecture cible} & \textbf{Point d'attention} \\
\midrule
Qlib \citep{githubQlib} & Plateforme Microsoft orientée recherche quantitative, données, modèles et backtesting. & Bon socle pour tester des facteurs, pipelines de données et modèles prédictifs. & Adaptation nécessaire aux données internes et aux contraintes de portefeuille réelles. \\
FinRL \citep{githubFinRL} et FinRL-Meta \citep{githubFinRLMeta} & Apprentissage par renforcement pour trading automatisé et environnements de marché. & Utile pour comparer une approche agentique LLM avec des agents RL plus classiques. & Le RL financier souffre fortement du surapprentissage et des coûts de transaction mal modélisés. \\
FinRL-X \citep{finrlx2026} & Infrastructure modulaire récente pour trading quantitatif natif IA. & Référence pour étudier l'intégration des briques de trading dans une infrastructure modulaire. & Projet récent : maturité opérationnelle à vérifier avant usage critique. \\
QuantConnect Lean \citep{githubLean} & Moteur open source de backtesting et trading algorithmique multi-actifs. & Brique robuste pour simuler et exécuter des stratégies dans un cadre contrôlé. & L'agent doit rester au-dessus du moteur, sans contourner les règles d'exécution. \\
OpenBB \citep{githubOpenBB} & Plateforme ouverte d'accès aux données et analyses financières. & Source pratique pour créer un agent de recherche et de veille. & La qualité dépend des fournisseurs de données configurés et de leurs licences. \\
Freqtrade \citep{githubFreqtrade} & Bot de trading crypto avec backtesting et hyperoptimisation. & Utile pour une étude sur crypto-actifs, où l'accès aux données et exchanges est simple. & Domaine crypto moins représentatif des contraintes actions institutionnelles. \\
\bottomrule
\end{tabular}
\end{adjustbox}
\end{table}

Qlib aurait répondu à un besoin de comparaison de modèles prédictifs ou de pipelines de facteurs. Mon expérience n'apprend pas de modèle et part de quatre ETF : j'ai conservé un moteur Python plus simple. Lean reste une référence pour une simulation multi-actifs plus complète. OpenBB a surtout été étudié pour l'agent de veille proposé ; le backtest disposait déjà de son historique. Une intégration d'actualités demanderait de documenter leur provenance, leur licence et leur disponibilité à la date de décision.

\section{Projets financiers et enseignements retenus}

\begin{table}[H]
\centering
\caption{Solutions GitHub orientées agents financiers.}
\label{tab:github-finance-agents}
\scriptsize
\begin{adjustbox}{width=\textwidth}
\begin{tabular}{p{0.18\textwidth}p{0.27\textwidth}p{0.31\textwidth}p{0.20\textwidth}}
\toprule
\textbf{Projet GitHub} & \textbf{Principe} & \textbf{Ce qu'il apporte au sujet} & \textbf{Limite pour un desk régulé} \\
\midrule
TradingAgents \citep{githubTradingAgents} & Framework multi-agents imitant une société de trading avec analystes spécialisés, trader et gestion du risque. & C'est l'un des dépôts les plus proches de l'idée d'équipe multi-agents appliquée au trading. & À analyser comme démonstrateur de recherche ; il faut ajouter validation des données, limites de risque, audit et contrôle humain. \\
FinRobot \citep{githubFinRobot} & Plateforme d'agents LLM pour applications financières, notamment recherche actions, analyse et valorisation. & Automatisation de tâches d'analyse financière et de rédaction de synthèses structurées. & Plus orienté analyse et reporting que génération robuste d'ordres multifactoriels. \\
FinMem \citep{finmem2023,githubFinMem} & Agent de trading fondé sur mémoire stratifiée et profils de décision. & Montre l'intérêt de la mémoire long terme pour contextualiser les décisions de trading. & Les performances expérimentales ne suffisent pas à prouver une robustesse exploitable en production. \\
AI Hedge Fund \citep{githubAiHedgeFund} & Démonstrateur open source d'un fonds simulé composé d'agents d'investissement. & Le dépôt illustre une répartition des tâches entre agents d'analyse et agent de décision. & Le dépôt se présente comme preuve de concept, pas comme outil de trading réel. \\
FinGPT \citep{fingpt2023,githubFinGPT} & Écosystème de modèles de langage financiers open source. & Modèles financiers à intégrer ou adapter pour des tâches d'analyse spécialisées. & Un modèle financier n'est pas une architecture de décision ; il faut ajouter outils, gouvernance et backtesting. \\
\bottomrule
\end{tabular}
\end{adjustbox}
\end{table}

TradingAgents et FinRobot ont nourri deux aspects du projet : la séparation des rôles d'analyse et de risque, puis les tâches de reporting proches de la note étudiée. FinMem soulève la question des synthèses passées : les conserver pourrait apporter du contexte, mais aussi transmettre leurs erreurs aux décisions suivantes. Cette mémoire reste hors du prototype. Aucun des cinq projets financiers n'a été exécuté ; leurs performances ne sont donc pas comparées à celles du backtest.

\section{Statut des briques et portée de la revue}

LangGraph est utilisé dans le workflow de restitution. Qlib, Lean et OpenBB sont des références étudiées, sans participation aux calculs finaux. TradingAgents et FinRobot ont inspiré les rôles proposés ; FinRL et Freqtrade relèvent d'autres périmètres expérimentaux. Pour examiner les dépôts, j'ai regardé l'état, les transitions, les permissions, la persistance, la documentation et les tests. La due diligence datée des huit dépôts retenus figure en annexe~\ref{annexe:github}.

La revue a conduit à un cas que je pouvais vérifier à partir des artefacts produits : confronter une réponse DeepSeek aux chiffres sources et à une note déterministe. L'extension mesure ensuite les conséquences de poids proposés sur le PnL. Les interactions libres entre agents et la mémoire ne font pas partie de ces expériences.

Les articles scientifiques étayent les concepts et les résultats ; les textes officiels éclairent la gouvernance. Les prépublications sont signalées comme telles. Les dépôts et documentations décrivent les fonctionnalités et les conditions d'intégration. Google Scholar, ScienceDirect et JSTOR ont servi au repérage ; les références renvoient autant que possible à l'article original, au DOI ou à une version légale déposée.
